%! AP Statistics — Unit 4, Lesson 4.12 — Group Activity KEY
\documentclass[10pt]{article}
\usepackage{apstats-article}
\usepackage{apstats-key}

\begin{document}

\pageheader{Unit 4, Lesson 4.12}{Group Activity KEY: Geometric vs.\ Binomial}
\namepartnerperiod

\vspace{0.1in}

\begin{scenariobox}[Directions]{sky}
\small
Work through all three tiers. Show all formula setup. Use complete sentences
when interpreting parameters.
\end{scenariobox}

\vspace{0.1in}

% ── TIER R ────────────────────────────────────────────────────────────────────
\begin{tcolorbox}[colback=goldbg, colframe=goldacc, boxrule=1pt, arc=2mm,
    title={\bfseries Tier R \quad Remediate: Identifying the Distribution},
    fonttitle=\small]
\small

For each scenario, (a) identify the distribution (binomial, geometric, or neither),
(b) state the key distinction, and (c) if geometric or binomial, state $X$ and the
parameters.

\vspace{10pt}
\textbf{Scenario 1:} A coin is flipped until tails appears.
$X =$ the number of flips.
\begin{enumerate}[label=(\alph*), leftmargin=2em, itemsep=6pt]
  \item Distribution: \ans{Geometric}
  \item Key distinction: \ansline{$n$ is not fixed --- flip until the first tails.}
  \item Parameters: \ans{$X \sim \text{Geom}(0.5)$, $p = 0.5$}
\end{enumerate}

\vspace{10pt}
\textbf{Scenario 2:} A coin is flipped 10 times.
$X =$ the number of tails.
\begin{enumerate}[label=(\alph*), leftmargin=2em, itemsep=6pt]
  \item Distribution: \ans{Binomial}
  \item Key distinction: \ansline{$n = 10$ is fixed in advance; $X$ counts tails.}
  \item Parameters: \ans{$X \sim B(10, 0.5)$, $n = 10$, $p = 0.5$}
\end{enumerate}

\vspace{10pt}
\textbf{Scenario 3:} A tetrahedral die (4 sides) is rolled. The die is fair.
$X =$ the number of the first roll.
\begin{enumerate}[label=(\alph*), leftmargin=2em, itemsep=6pt]
  \item Distribution: \ans{Neither}
  \item Key distinction: \ansline{One roll only, values 1--4 each with prob.\ $1/4$ (uniform).}
\end{enumerate}
\vspace{4pt}
\end{tcolorbox}

\vspace{0.1in}

% ── TIER A ────────────────────────────────────────────────────────────────────
\begin{tcolorbox}[colback=sky, colframe=navy, boxrule=1pt, arc=2mm,
    title={\bfseries Tier A \quad Approaching Proficiency: Full Geometric Analysis},
    fonttitle=\small]
\small

\textbf{Context:} An archer hits the bullseye 30\% of the time. She shoots arrows
one at a time until hitting the bullseye.
Let $X =$ the number of arrows shot until (and including) the first bullseye.

\vspace{8pt}
(a) Check BITS. Write one sentence per condition.

\vspace{4pt}
B: \ansline{Each arrow hits the bullseye (success) or misses (failure).}

I: \ansline{One shot does not affect the next.}

T: \ansline{She shoots until the first bullseye.}

S: \ansline{Same $p = 0.30$ on every shot.}

\vspace{6pt}
State the distribution: $X \sim $ \ans{Geom(0.30)}

\vspace{8pt}
(b) Compute $P(X = 5)$. Show the formula and the numerical answer.
\[
  P(X = 5) = ({\color{keyred}\mathbf{0.70}})^{4}({\color{keyred}\mathbf{0.30}}) = {\color{keyred}\mathbf{(0.2401)(0.30)}} \hspace{2cm} \approx {\color{keyred}\mathbf{0.0720}}
\]

\vspace{8pt}
(c) Compute $\mu_X$ and interpret in context.
\vspace{4pt}
$\mu_X = $ \ans{$1/p = 1/0.30$} $= $ \ans{3.33}

\vspace{4pt}
Interpretation: \ansline{Over many rounds, she averages}\\[1pt]\ansline{about 3.33 arrows to get her first bullseye.}

\vspace{8pt}
(d) Find $P(X \le 4)$ using \texttt{geometcdf}. What does this represent?
\vspace{4pt}
$P(X \le 4) = \texttt{geometcdf(}{\color{keyred}\mathbf{0.30}}\texttt{,}\;4\texttt{)} \approx $ \ans{0.7599}

\vspace{4pt}
Interpretation: \ansline{About a 76\% chance she hits}\\[1pt]\ansline{her first bullseye within her first 4 arrows.}
\vspace{4pt}
\end{tcolorbox}

\vspace{0.1in}

% ── TIER E ────────────────────────────────────────────────────────────────────
\begin{tcolorbox}[colback=greenbg, colframe=greenacc, boxrule=1pt, arc=2mm,
    title={\bfseries Tier E \quad Extension: Distribution Table and Shape},
    fonttitle=\small]
\small

\textbf{Context (same archer):} $X \sim \text{Geom}(0.30)$.

\vspace{8pt}
(a) Build the partial distribution table for $X = 1, 2, 3, 4, 5$.

\vspace{4pt}
\renewcommand{\arraystretch}{1.8}
\begin{tabularx}{\linewidth}{@{} >{\bfseries}p{1.2cm} X X X X X @{}}
$k$ & 1 & 2 & 3 & 4 & 5 \\
\hline
$P(X=k)$ & \ans{0.3000} & \ans{0.2100} & \ans{0.1470} & \ans{0.1029} & \ans{0.0720} \\
\end{tabularx}

\vspace{8pt}
(b) Describe the shape of the geometric distribution based on your table.
Does it look like the binomial distribution for the same $p$?
\vspace{4pt}
\ansline{Strictly decreasing --- strongly right-skewed, peak at $k = 1$.}\\[1pt]\ansline{No: a binomial with the same $p$ is mound-shaped near $np$.}

\vspace{8pt}
(c) Now consider $Y \sim B(5, 0.30)$ (same $p$, fixed $n = 5$).
Compare $P(X = 5)$ and $P(Y = 5)$. Explain why they are different,
and state what each probability means in context.
\vspace{4pt}
\ansline{$P(X=5) \approx 0.0720$: 1st bullseye on arrow 5.}\\[1pt]\ansline{$P(Y=5) = (0.30)^5 \approx 0.0024$: all 5 of 5 arrows hit the bullseye.}\\[1pt]\ansline{Different events: $X=5$ is 4 misses then a hit; $Y=5$ is 5 hits.}
\vspace{4pt}
\end{tcolorbox}

\end{document}
